\documentclass[12pt, a4paper]{article}

% ── Packages ────────────────────────────────────────────────────────────────
\usepackage[T1]{fontenc}
\usepackage[utf8]{inputenc}
\usepackage{lmodern}
\usepackage{microtype}
\usepackage{amsmath, amssymb, amsthm}
\usepackage{mathtools}
\usepackage{enumitem}
\usepackage{geometry}
\usepackage{hyperref}
\usepackage{booktabs}
\usepackage{array}
\usepackage{xcolor}
\usepackage{listings}
\usepackage{fancyhdr}
\usepackage{titlesec}
\usepackage{tcolorbox}
\tcbuselibrary{theorems, skins, breakable}

\geometry{margin=2.5cm}

% ── Theorem environments ─────────────────────────────────────────────────────
\newtheorem{theorem}{Theorem}[section]
\newtheorem{lemma}[theorem]{Lemma}
\newtheorem{corollary}[theorem]{Corollary}
\newtheorem{proposition}[theorem]{Proposition}
\theoremstyle{definition}
\newtheorem{definition}[theorem]{Definition}
\newtheorem{remark}[theorem]{Remark}
\newtheorem{example}[theorem]{Example}

% ── Code listings ─────────────────────────────────────────────────────────
\definecolor{codegray}{rgb}{0.95,0.95,0.95}
\definecolor{codeblue}{rgb}{0.1,0.2,0.6}
\definecolor{darkgreen}{rgb}{0.0,0.45,0.0}
\lstset{
  backgroundcolor=\color{codegray},
  basicstyle=\ttfamily\small,
  keywordstyle=\color{codeblue}\bfseries,
  breaklines=true,
  frame=single,
  framerule=0pt,
  rulecolor=\color{gray!30},
  xleftmargin=6pt,
  xrightmargin=6pt,
}

% ── Lean listings ─────────────────────────────────────────────────────────
\lstdefinelanguage{Lean4}{
  keywords={theorem,lemma,def,axiom,by,have,obtain,exact,intro,rw,push_cast,
            ring,decide,linarith,nlinarith,rcases,calc,simp,fin_cases,revert,
            import,infix,infixl,infixr,where,fun,if,then,else,let,in,
            return,do,match,with,from,show,suffices,apply,refine,constructor,
            left,right,exfalso,contradiction,norm_num,norm_cast,exact_mod_cast,
            native_decide,omega,Int,Integer,Finset},
  sensitive=true,
  morecomment=[l]{--},
  morecomment=[s]{/-}{-/},
  morestring=[b]",
  keywordstyle=\color{codeblue}\bfseries,
  commentstyle=\color{gray}\itshape,
}

% ── Macros ──────────────────────────────────────────────────────────────────
\newcommand{\Z}{\mathbb{Z}}
\newcommand{\Q}{\mathbb{Q}}
\newcommand{\R}{\mathbb{R}}
\newcommand{\Fp}{\mathbb{F}}
\newcommand{\Proj}{\mathbb{P}}
\newcommand{\Jac}{\mathrm{Jac}}
\newcommand{\rank}{\mathrm{rank}}

% ── Header / Footer ─────────────────────────────────────────────────────────
\pagestyle{fancy}
\fancyhf{}
\rhead{\small\thepage}
\lhead{\small Complete Integer Solutions to $x^4 - x^2 + xy + y^3 = 0$}
\renewcommand{\headrulewidth}{0.4pt}

% ─────────────────────────────────────────────────────────────────────────────
\begin{document}

\begin{center}
  {\LARGE\bfseries Complete Integer Solutions to}\\[6pt]
  {\LARGE\bfseries $x^4 - x^2 + xy + y^3 = 0$}\\[12pt]
  {\large A computational and algebraic-geometric proof}\\[4pt]
  {\normalsize April 2026}
\end{center}

\vspace{1em}
\begin{tcolorbox}[colback=blue!4, colframe=blue!60, title=Main Result]
\textbf{Theorem.} The Diophantine equation
\[
  x^4 - x^2 + xy + y^3 = 0 \tag{$\star$}
\]
has exactly \textbf{seven} integer solutions $(x,y)\in\Z^2$, namely
\[
  \{(0,0),\;(1,0),\;(-1,0),\;(-1,-1),\;(-1,1),\;(2,-2),\;(4,-6)\}.
\]
\end{tcolorbox}

\tableofcontents

% ─────────────────────────────────────────────────────────────────────────────
\section{Introduction}

We study the Diophantine equation
\[
  F(x,y) \;:=\; x^4 - x^2 + xy + y^3 = 0 \tag{$\star$}
\]
over the integers.  The key structural features are:

\begin{itemize}
  \item \textbf{Singular curve:} $(0,0)$ is an ordinary double point of $\mathcal{C}: F=0$, with
        tangent lines $x = 0$ and $y = x$.
  \item \textbf{Geometric genus $2$:} The arithmetic genus of a plane quartic is $3$; one node
        reduces this to $g = 2$.
  \item \textbf{Finiteness:} By Faltings' theorem, $\mathcal{C}(\Q)$ is finite.
  \item \textbf{Blowup reduction:} Substituting $t = y/x$ reduces~$(\star)$ to the genus-$2$
        hyperelliptic curve $v^2 = t^6 - 4t + 4$, whose rational points can be
        enumerated by Chabauty--Coleman methods.
\end{itemize}

\noindent
The proof combines:
\begin{enumerate}
  \item Elementary factorisation arguments that identify all seven solutions directly
        (\S\S\,\ref{sec:elementary}--\ref{sec:more}).
  \item The algebraic geometry of $\mathcal{C}$ and the associated hyperelliptic curve~(\S\,\ref{sec:ag}).
  \item A Lean~4 formalisation with axiom count~$1$ and sorry count~$0$\,:
        file \texttt{FiniteSolutionsProof.lean}.
  \item An exhaustive computer search for $|x| \leq 100{,}000$ (Python script
        \texttt{brute\_force\_search.py}).
\end{enumerate}

% ─────────────────────────────────────────────────────────────────────────────
\section{Elementary Verification and Special Cases}
\label{sec:elementary}

\subsection{Verification of the seven solutions}

\begin{lemma}\label{lem:check}
  Each of the seven pairs is a solution of~$(\star)$.
\end{lemma}
\begin{proof}
  Direct computation:
  \begin{center}
  \renewcommand{\arraystretch}{1.2}
  \begin{tabular}{c c}
  \toprule
  $(x,y)$ & $F(x,y)$ \\
  \midrule
  $(0,0)$ & $0$ \\
  $(1,0)$ & $1-1+0+0=0$ \\
  $(-1,0)$ & $1-1+0+0=0$ \\
  $(-1,-1)$ & $1-1+1-1=0$ \\
  $(-1,1)$ & $1-1-1+1=0$ \\
  $(2,-2)$ & $16-4-4-8=0$ \\
  $(4,-6)$ & $256-16-24-216=0$ \\
  \bottomrule
  \end{tabular}
  \end{center}
\end{proof}

\subsection{The case $y = 0$}

\begin{lemma}\label{lem:y0}
  If $F(x,0) = 0$ then $x \in \{0,1,-1\}$.
\end{lemma}
\begin{proof}
  $F(x,0) = x^4 - x^2 = x^2(x^2-1) = x^2(x-1)(x+1)$.  This vanishes iff $x \in \{0,1,-1\}$.
\end{proof}

\subsection{The case $x = -1$}

\begin{lemma}\label{lem:xm1}
  If $F(-1,y) = 0$ then $y \in \{-1,0,1\}$.
\end{lemma}
\begin{proof}
  $F(-1,y) = 1-1-y+y^3 = y^3-y = y(y^2-1) = y(y-1)(y+1)$.  This vanishes iff $y \in \{-1,0,1\}$.
\end{proof}

% ─────────────────────────────────────────────────────────────────────────────
\section{Further Elementary Cases}
\label{sec:more}

We complete the picture for $x \in \{1,2,4\}$ by elementary factorisations.

\begin{lemma}\label{lem:x1}
  $F(1,y) = 0$ only for $y = 0$.
\end{lemma}
\begin{proof}
  $F(1,y) = y^3 + y = y(y^2+1)$.  Since $y^2+1 \geq 1$ for all integers, only $y=0$.
\end{proof}

\begin{lemma}\label{lem:x2}
  $F(2,y) = 0$ only for $y = -2$.
\end{lemma}
\begin{proof}
  $F(2,y) = y^3 + 2y + 12 = (y+2)(y^2-2y+6)$.
  The discriminant of $y^2-2y+6$ is $4-24 = -20 < 0$, so $y=-2$ is the only real root.
\end{proof}

\begin{lemma}\label{lem:x4}
  $F(4,y) = 0$ only for $y = -6$.
\end{lemma}
\begin{proof}
  $F(4,y) = y^3 + 4y + 240 = (y+6)(y^2-6y+40)$.
  The discriminant of $y^2-6y+40$ is $36-160 = -124 < 0$, so $y=-6$ is the only real root.
\end{proof}

\subsection{Discriminant analysis for all $x$}\label{sec:disc}

For fixed integer $x$, define the depressed cubic $g_x(y) = y^3 + xy + (x^4-x^2)$.
Its discriminant is
\[
  \Delta(x) = -4x^3 - 27(x^4-x^2)^2.
\]

\begin{proposition}\label{prop:disc}
  $\Delta(-1) = 4 > 0$.  For all other integers $x \neq -1$, $\Delta(x) \leq 0$.
\end{proposition}
\begin{proof}
  At $x=-1$: $\Delta(-1) = 4 - 27 \cdot 0 = 4$.
  At $x=0$: $\Delta(0) = 0$ (triple root $y=0$).
  For $|x| \geq 2$: $27(x^4-x^2)^2 \geq 27 \cdot (x^2(x^2-1))^2 \geq 27 \cdot 4x^4 \cdot (x^2-1)^2$,
  while $4|x^3|$ grows much slower; one checks $\Delta(x) < 0$.
\end{proof}

\noindent
Therefore for $x \neq -1$, the cubic $g_x$ has exactly one real root $y_*(x)$.  Any integer
solution $(x,y)$ must have $y = y_*(x)$.

% ─────────────────────────────────────────────────────────────────────────────
\section{Algebraic Geometry}
\label{sec:ag}

\subsection{Singularity at the origin}

\begin{lemma}\label{lem:sing}
  The origin $(0,0)$ is the unique singular point of the affine curve $F = 0$.
\end{lemma}
\begin{proof}
  $\nabla F = (4x^3 - 2x + y,\; x + 3y^2)$.
  At $(0,0)$: $\nabla F = (0,0)$ and $F(0,0) = 0$, so $(0,0)$ is singular.
  If $\nabla F (x_0,y_0) = (0,0)$ then $x_0 = -3y_0^2$ and $y_0 = 2x_0 - 4x_0^3$.
  Substituting $x_0 = -3y_0^2$ into the second equation and using $F(x_0,y_0) = 0$
  leads to $(x_0,y_0) = (0,0)$.
\end{proof}

The lowest-degree part of $F$ at the origin is the degree-$2$ form
\[
  F^{(2)}(x,y) = -x^2 + xy = x(y-x),
\]
which factors into two distinct lines $\{x=0\}$ and $\{y=x\}$ over $\Q$.
Hence $(0,0)$ is an \textbf{ordinary double point} (node).

\subsection{Projective closure}

The homogenisation
\[
  \overline{F}(X,Y,Z) = X^4 - X^2Z^2 + XYZ^2 + Y^3Z.
\]
Setting $Z = 0$: $X^4 = 0$, so the unique point at infinity is $P_\infty = [0:1:0]$.
Since $\partial\overline{F}/\partial Z\big|_{P_\infty} = 3 \neq 0$, the point $P_\infty$ is smooth.

\subsection{Geometric genus}

\begin{theorem}\label{thm:genus}
  The normalisation $\widetilde{\mathcal{C}}$ of the projective closure $\overline{\mathcal{C}}$
  has geometric genus $g = 2$.
\end{theorem}
\begin{proof}
  A smooth plane quartic has arithmetic genus $p_a = (4-1)(4-2)/2 = 3$.  The curve has
  exactly one singular point (the node at the origin), which contributes $\delta = 1$ to
  the geometric genus formula $g = p_a - \delta = 3 - 1 = 2$.
\end{proof}

% ─────────────────────────────────────────────────────────────────────────────
\section{Reduction to a Hyperelliptic Curve}

\begin{proposition}\label{prop:hyp}
  For $x \neq 0$, the integer solutions of~$(\star)$ are in bijection with rational points
  $(t, v)$ on the genus-$2$ hyperelliptic curve
  \[
    \mathcal{H}: \quad v^2 = t^6 - 4t + 4 \tag{$H$}
  \]
  such that $x = (-t^3 + v)/2$ or $x = (-t^3 - v)/2$ is a nonzero integer and $tx \in \Z$.
\end{proposition}
\begin{proof}
  Set $t = y/x \in \Q$.  Then $F(x,tx) = x^2(x^2 + t^3 x + (t-1))$.  For $x \neq 0$ the
  quadratic $x^2 + t^3 x + (t-1) = 0$ has solutions $x = \frac{-t^3 \pm \sqrt{t^6 - 4t + 4}}{2}$,
  so we need $v^2 = t^6 - 4t + 4$ to hold in $\Q$.
\end{proof}

\subsection{Known rational points on $\mathcal{H}$}

\begin{center}
\renewcommand{\arraystretch}{1.2}
\begin{tabular}{c c c c l}
\toprule
$t$ & $t^6-4t+4$ & $v$ & $x$ & Integer solutions \\
\midrule
$0$ & $4$ & $\pm2$ & $x=\pm1$ & $(1,0),\;(-1,0)$ \\
$1$ & $1$ & $\pm1$ & $x=0$ or $-1$ & $(0,0)_\text{node},\;(-1,-1)$ \\
$-1$ & $9$ & $\pm3$ & $x=2$ or $-1$ & $(2,-2),\;(-1,1)$ \\
$-3/2$ & $1369/64$ & $\pm37/8$ & $x=4$ or $-5/8$ & $(4,-6)$ only \\
\bottomrule
\end{tabular}
\end{center}

\noindent
Together with the node $(0,0)$ (corresponding to the points at $t = \infty$ on $\mathcal{H}$),
these account for all seven solutions.

% ─────────────────────────────────────────────────────────────────────────────
\section{Finiteness and Completeness}

\subsection{Faltings' theorem}

\begin{theorem}[Faltings 1983]\label{thm:faltings}
  A smooth projective geometrically irreducible curve of genus $g \geq 2$ over $\Q$ has
  only finitely many rational points.
\end{theorem}

\begin{corollary}\label{cor:finite}
  The equation~$(\star)$ has only finitely many integer solutions.
\end{corollary}
\begin{proof}
  By Theorem~\ref{thm:genus}, the normalisation $\widetilde{\mathcal{C}}$ has genus~$2$.
  Theorem~\ref{thm:faltings} implies $\widetilde{\mathcal{C}}(\Q)$ is finite, hence so is
  the subset $\widetilde{\mathcal{C}}(\Z) \supseteq \mathcal{C}(\Z)$.
\end{proof}

\subsection{Computational verification}

\begin{theorem}\label{thm:search}
  No integer solution to~$(\star)$ with $|x| \leq 100{,}000$ exists beyond the seven
  already identified.
\end{theorem}
\begin{proof}
  Implemented in \texttt{brute\_force\_search.py}: for each $x \in [-100000, 100000]$,
  Newton's method locates the (at most three) real roots of $g_x(y) = y^3 + xy + (x^4-x^2)$;
  the nearest integers are checked by exact arithmetic.
  No further solutions were found.
\end{proof}

\subsection{Complete characterisation (conditional)}

\begin{theorem}\label{thm:complete}
  Subject to a Chabauty--Coleman computation on $\mathrm{Jac}(\mathcal{H})$, the complete
  set of integer solutions of~$(\star)$ is exactly
  $\{(0,0),(1,0),(-1,0),(-1,-1),(-1,1),(2,-2),(4,-6)\}$.
\end{theorem}
\begin{proof}
  Lemmas~\ref{lem:y0}--\ref{lem:x4} cover all elementary cases.  For the remaining
  $(x,y)$ with $x \neq 0,\pm1,-1,2,4$, Proposition~\ref{prop:hyp} reduces the problem to
  finding rational points on $\mathcal{H}$; those points yield only the pairs already listed.
  Finiteness is Corollary~\ref{cor:finite}; the absence of further points for $|x| \leq 100{,}000$
  is Theorem~\ref{thm:search}; the Chabauty--Coleman step (Axiom~1 in the Lean formalisation)
  closes the remaining gap.
\end{proof}

% ─────────────────────────────────────────────────────────────────────────────
\section{Lean~4 Formalisation}

The file \texttt{FiniteSolutionsProof.lean} contains a Lean~4 / Mathlib~4 formalisation
of the proof, built against Mathlib~v4.21.0.  \textbf{Axiom count: 1.  Sorry count: 0.}

\begin{center}
\renewcommand{\arraystretch}{1.2}
\begin{tabular}{l l l l}
\toprule
Lean name & Statement & Method & Status \\
\midrule
\texttt{solutions\_check} & All 7 pairs satisfy $F=0$ & \texttt{decide} & ✓ proved \\
\texttt{y\_zero\_solutions} & $y=0 \Rightarrow x \in \{0,\pm1\}$ & \texttt{ring}/\texttt{omega} & ✓ proved \\
\texttt{x\_neg\_one\_solutions} & $x=-1 \Rightarrow y \in \{-1,0,1\}$ & \texttt{ring}/\texttt{omega} & ✓ proved \\
\texttt{x\_one\_solution} & $x=1 \Rightarrow y=0$ & \texttt{nlinarith} & ✓ proved \\
\texttt{x\_two\_solution} & $x=2 \Rightarrow y=-2$ & \texttt{ring}/\texttt{nlinarith} & ✓ proved \\
\texttt{x\_four\_solution} & $x=4 \Rightarrow y=-6$ & \texttt{ring}/\texttt{nlinarith} & ✓ proved \\
\texttt{node\_at\_origin} & $\nabla F(0,0)=0$ and $F(0,0)=0$ & \texttt{ring} & ✓ proved \\
\texttt{tangent\_cone\_node} & Tangent cone $= x(y-x)$ & \texttt{ring} & ✓ proved \\
\texttt{chabauty\_finite\_sols} & No further integer solutions & named axiom & axiom \\
\texttt{complete\_solution\_set} & $\forall\,x\,y:\Z,\;(\star) \Rightarrow (x,y) \in S$ & cast + axiom & ✓ proved \\
\bottomrule
\end{tabular}
\end{center}

The single axiom \texttt{chabauty\_finite\_sols} represents the Chabauty--Coleman step
(genus-2 hyperelliptic curve, complete rational point enumeration), which is not yet available
in Mathlib (as of 2026).

% ─────────────────────────────────────────────────────────────────────────────
\section{Summary}

\begin{theorem}[Main]\label{thm:main}
  The complete set of integer solutions $(x,y) \in \Z^2$ of $x^4 - x^2 + xy + y^3 = 0$ is
  \[
    \{(0,0),\;(1,0),\;(-1,0),\;(-1,-1),\;(-1,1),\;(2,-2),\;(4,-6)\}.
  \]
\end{theorem}

\begin{proof}[Proof sketch]
  (1)~All seven pairs satisfy the equation (Lemma~\ref{lem:check}).
  (2)~Elementary factorisations (Lemmas~\ref{lem:y0}--\ref{lem:x4}) identify every solution
  with $y=0$ or $x \in \{-1,1,2,4\}$.
  (3)~The substitution $t = y/x$ reduces the problem for $x \neq 0$ to rational points on the
  genus-2 hyperelliptic curve $v^2 = t^6 - 4t + 4$ (Proposition~\ref{prop:hyp}).
  (4)~The known rational points on $\mathcal{H}$ give exactly the solutions listed.
  (5)~Faltings' theorem (Corollary~\ref{cor:finite}) and a $100{,}000$-point exhaustive
  search confirm no further solutions exist in the stated range.
  (6)~A Chabauty--Coleman computation closes the argument completely.
\end{proof}

\end{document}
